\documentclass{ccnudoc}
\usepackage{multirow,xeCJKfntef,xpinyin}
\hypersetup{
  pdftitle  = {CCNUthesis: 华中师范大学学位论文模版},
  pdfauthor = {夏康玮}}
% 全角标点放在引号中，需要改成半角式，否则间距过大，不好看
\def\FSID{“{\xeCJKsetup{PunctStyle=banjiao}。}”} % U+3002
\def\FSFW{“{\xeCJKsetup{PunctStyle=banjiao}．}”} % U+FF0E
\def\COFW{“{\xeCJKsetup{PunctStyle=banjiao}：}”} % U+FF1A
\def\SCFW{“{\xeCJKsetup{PunctStyle=banjiao}；}”} % U+FF1B


\title{\textcolor{MaterialIndigo800}{%
  \textbf{CCNUthesis: 华中师范大学学位论文\xpinyin[font=\sffamily,format=\color{MaterialIndigo800}]{模}{mu2}板}}}
\author{夏康玮}
\date{2022/05/22\quad v1.2.1%
  \thanks{\url{https://github.com/xkwxdyy/CCNUthesis}.}}


\begin{document}
\newgeometry{
  left   = 1.25 in,
  right  = 1.25 in,
  top    = 1.25 in,
  bottom = 1.00 in
}

\maketitle
\vfill
\begin{tikzpicture}[remember picture, overlay]
  \node[opacity = 0.2] at ([shift={(0.1\textwidth, -0.45\textheight)}]current page text area.north west){
    \includegraphics[width=26cm]{ccnu-logo.png}
  };
\end{tikzpicture}
\vfill
\thispagestyle{plain}
\clearpage

\tableofcontents


% 用户手册的页边距

\newgeometry{
  left   = 1.65 in,
  right  = 0.80 in,
  top    = 1.25 in,
  bottom = 1.00 in
}


\section{还有什么没写：TODO}

\begin{itemize}
  \item 减少宏包依赖
  \item 如果是 author-year 的时候，对于多音字怎么办？——加 key
  \item 完整例子、physicx、tabularray示例和用户使用的分开
  \item 重写 Readme
  \item 重写 wiki
  \item 把 manual 中有用的部分写进 wiki
  \item 写用户手册
    \begin{itemize}
      \item 如何提问
    \end{itemize}
  \item 看 AlphaZTX 的知乎
  \item 学习其它论文模版的处理
    \begin{itemize}
      \item thuthesis
      \item sjuthesis
      \item njuthesis
      \item hithesis
    \end{itemize}
  \item 处理硕士模版
\end{itemize}

\section{介绍}

目前，在网上可以找到的复旦大学 \LaTeX{} 论文模板主要有以下这些：
\begin{itemize}
  \item 数学科学学院 2001 级的何力同学和李湛同学在 2005 年根据
    学校要求所设计的 \cls{毕业论文格式 tex04 版}，以及 2008 年
    张越同学修改之后的 \cls{毕业论文格式 tex08 版}，这是专为
    数院本科生撰写毕业论文而设计的
    \cite{数院毕业论文格式,数院毕业论文格式更新}；
  \item Pandoxie 编写的 \cls{FDU-Thesis-Latex}
    \cite{pandoxie2014fduthesislatex}，基本满足了博士（硕士）
    毕业论文格式要求，使用人数较多；
  \item richarddzh 编写的硕士论文模板 \cls{fudan-thesis}
    \cite{richard2016fudanthesis}。
\end{itemize}
以上这些模板大都没有经过系统的设计，也鲜有后续维护。相比之下，
清华大学 \cite{thuthesis}、重庆大学 \cite{cquthesis}、
中国科学技术大学 \cite{ustcthesis}、中国科学院大学 \cite{ucasthesis}
以及友校上海交通大学 \cite{sjtuthesis}等，都有成熟、
稳定的解决方案，值得参考。

本模板将借鉴前辈经验，重新设计，并使用 \LaTeX3
\cite{source3} 编写，以适应 \TeX{} 技术发展潮流；
同时还将构建一套简洁的接口，方便用户使用。

\subsection*{\LaTeX{} 入门}

本文档并非是一份 \LaTeX{} 零基础教程。如果您是完完全全的新手，
建议先阅读相关入门文档，如刘海洋编著的《\LaTeX{} 入门》
\cite{刘海洋2013latex入门} 第一章，或大名鼎鼎的“\pkg{lshort}”
\cite{lshort} 及其中文翻译版 \cite{lshort-zh-cn}。当然，
网络上的入门教程多如牛毛，您可以自行选取。

\subsection*{关于本文档}

本文采用不同字体表示不同内容。无衬线字体表示宏包名称，如
\pkg{xeCJK} 宏包、\cls{fduthesis} 文档类等；等宽字体表示代码或
文件名，如 \cs{fdusetup} 命令、\env{abstract} 环境、\TeX{} 文档
\file{thesis.tex} 等；带有尖括号的楷体（或西文斜体）表示命令参数，
如 \meta{模板选项}、\meta{English title} 等。在使用时，参数两侧
的尖括号不必输入。示例代码进行了语法高亮处理，以方便阅读。

在用户手册中，带有蓝色侧边线的为 \LaTeX{} 代码，而带有粉色侧边线
的则为命令行代码，请注意区分。模板提供的选项、命令、环境等，
均用横线框起，同时给出使用语法和相关说明。

本模板中的选项、命令或环境可以分为以下三类：
\begin{itemize}
  \item 名字后面带有 \rexptarget\rexpstar{} 的，表示只能
    在\emph{中文模板}中使用；
  \item 名字后面带有 \exptarget\expstar{} 的，表示只能
    在\emph{英文模板}中使用；
  \item 名字后面不带有特殊符号的，表示既可以在中文模板中使用，
    也可以在英文模板中使用。
\end{itemize}

代码实现主要面向对 \LaTeX{} 宏包开发感兴趣的用户。如果您有任何改进
意见或者功能需求，欢迎前往 GitHub 仓库
\href{https://github.com/stone-zeng/fduthesis/issues}{提交 issue}。

文档的最后还提供了版本历史和代码索引，以供查阅。

\section{安装}

\subsection{获取 \cls{fduthesis}}

\subsubsection{标准安装}

如果没有特殊理由，始终建议您使用宏包管理器安装 \cls{fduthesis}。
例如在 \TeXLive{} 中，执行（可能需要管理员权限）
\begin{shellexample}[morekeywords={tlmgr,install}]
  tlmgr install fduthesis
\end{shellexample}
即可完成安装。

在 \TeXLive{} 和 \MiKTeX{} 中，您还可以通过图形界面进行安装，
此处不再赘述。

\subsubsection{手动安装}

如果您需要从 CTAN 上自行下载并手动安装，较好的方法是使用 TDS
安装包：
\begin{itemize}
  \item 从 CTAN 上下载 \cls{fduthesis} 的
    \href{http://mirror.ctan.org/install/macros/latex/contrib/fduthesis.tds.zip}{TDS 安装包}；
  \item 按目录结构将 \file{fduthesis.tds.zip} 中的文件复制到 \TeX{}
    发行版的本地 TDS 根目录；
  \item 执行 \bashcmd{mktexlsr} 刷新文件名数据库以完成安装。
\end{itemize}
%
您也可以从源代码直接生成模板（不推荐）：
\begin{itemize}
  \item 打开 \href{https://github.com/stone-zeng/fduthesis}{项目主页}，
    点击“Code”按钮，并选择“Download ZIP”，下载 \file{fduthesis-main.zip}；
    如果您的电脑中安装有 git 程序，也可通过以下命令直接克隆代码仓库：
    \begin{shellexample}[gobble=7,alsoletter={.},morekeywords={git,clone}]
      git clone https://github.com/stone-zeng/fduthesis.git
    \end{shellexample}
  \item 解压并进入到 \file{source} 文件夹，执行以下命令以生成
    模板的各组件：
    \begin{shellexample}[gobble=7,morekeywords={xetex}]
      xetex fduthesis.dtx
    \end{shellexample}
  \item 将生成的文档类（\file{.cls}）、宏包（\file{.sty}）以及
    参数配置文件（\file{.def}）复制到 \TeX{} 发行版本地 TDS 树
    的 \path{texmf-local/tex/latex/fduthesis/} 目录下，并执行
    \bashcmd{mktexlsr} 刷新文件名数据库，方可完成安装。
  \item 使用 \cls{fduthesis} 撰写论文时，您还需要从代码仓库下的
    \file{testfiles/support} 目录中复制 \file{fudan-name.pdf}
    文件至工作目录，以确保封面中的校名图片可以正确显示。
\end{itemize}

\subsubsection{扁平化安装}

如果您不希望安装本模板，但需要立刻使用，也可以使用模板提供的安装脚本。
从 GitHub 上获取代码仓库后，执行 \file{install-win.bat}（Windows 系统）
或 \file{install-linux.sh}（Linux 系统），所有需要的文件便会在
\file{thesis} 文件夹中生成。

\subsection{模板组成}

本模板主要包含核心文档类、配置文件、附属宏包以及用户文档等几个部分，
其具体组成见表~\ref{tab:fduthesis-components}。

\begin{table}[ht]
  \caption{\cls{fduthesis} 的主要组成部分}
  \label{tab:fduthesis-components}
  \centering
  \small
  \begin{tabular}{ll}
    \toprule
      \textbf{文件} & \textbf{功能说明} \\
    \midrule
      \file{fduthesis.cls}          & 中文模板文档类 \\
      \file{fduthesis-en.cls}       & 英文模板文档类 \\
      \file{fduthesis.def}          & 参数配置文件，用于设定
        \cls{fduthesis} 的初始参数，不建议您自行改动 \\
      \file{fdudoc.cls}             & 用户手册文档类 \\
      \file{fdulogo.sty}            & 复旦大学视觉识别系统 \\
      \file{fudan-emblem.pdf}       & 校徽 \\
      \file{fudan-emblem-new.pdf}   & 校徽（重修版） \\
      \file{fudan-name.pdf}         & 校名图片 \\
      \file{README.md}              & 简要自述 \\
      \ifdefined\FDUCODEDOC
        \file{fduthesis.pdf}        & 中文用户手册 \\
        \file{fduthesis-en.pdf}     & 英文用户手册 \\
        \file{fduthesis-code.pdf}   & 模板实现代码（本文档） \\
      \else
        \file{fduthesis.pdf}        & 中文用户手册（本文档） \\
        \file{fduthesis-en.pdf}     & 英文用户手册 \\
        \file{fduthesis-code.pdf}   & 模板实现代码 \\
      \fi
      \file{fduthesis-template.tex} & 空白模板，可据此为基础完成论文
        撰写 \\
    \bottomrule
  \end{tabular}
\end{table}

\section{使用说明}

\subsection{基本用法}

以下是一份简单的 \TeX{} 文档，它演示了 \cls{fduthesis} 的最基本用法：

\begin{latexexample}[deletetexcs={\documentclass},
    moretexcs={\chapter},morekeywords={\documentclass},
    emph={[2]document}]
  % thesis.tex
  \documentclass{fduthesis}
  \begin{document}
    \chapter{欢迎}
    \section{Welcome to fduthesis!}
    你好，\LaTeX{}！
  \end{document}
\end{latexexample}


按照 \ref{subsec:编译方式}~小节中的方式编译该文档，您应当得到
一篇 5 页的文章。当然，这篇文章的绝大部分都是空白的。

英文模板可以用类似的方式使用：

\begin{latexexample}[deletetexcs={\documentclass},
    moretexcs={\chapter},morekeywords={\documentclass},
    emph={[2]document}]
  % thesis-en.tex
  \documentclass{fduthesis-en}
  \begin{document}
    \chapter{Welcome}
    \section{Welcome to fduthesis!}
    Hello, \LaTeX{}!
  \end{document}
\end{latexexample}

英文模板只对正文部分进行了改动，封面、指导小组成员以及声明页仍将
显示为中文。

\subsection{编译方式} \label{subsec:编译方式}

本模板不支持 \pdfTeX{} 引擎，请使用 \XeLaTeX{} 或 \LuaLaTeX{}
编译。推荐使用 \XeLaTeX{}。为了生成正确的目录、脚注以及交叉引用，
您至少需要连续编译两次。

以下代码中，假设您的 \TeX{} 源文件名为 \file{thesis.tex}。
使用 \XeLaTeX{} 编译论文，请在命令行中执行
\begin{shellexample}[morekeywords={xelatex}]
  xelatex thesis
\end{shellexample}
或使用 \pkg{latexmk}：
\begin{shellexample}[morekeywords={latexmk},emph={-xelatex}]
  latexmk -xelatex thesis
\end{shellexample}

使用 \LuaLaTeX{} 编译论文，请在命令行中执行
\begin{shellexample}[morekeywords={lualatex}]
  lualatex thesis
\end{shellexample}
或者
\begin{shellexample}[morekeywords={latexmk},emph={-lualatex}]
  latexmk -lualatex thesis
\end{shellexample}

\subsection{模板选项}

所谓“模板选项”，指需要在引入文档类的时候指定的选项：
\begin{latexexample}[deletetexcs={\documentclass},
    morekeywords={\documentclass}]
  \documentclass(*\oarg{模板选项}*){fduthesis}
\end{latexexample}

有些模板选项为布尔型，它们只能在 \opt{true} 和 \opt{false}
中取值。对于这些选项，\kvopt{\meta{选项}}{true} 中的“|= true|”
可以省略。


\begin{function}[added=2018-02-01]{type}
  \begin{fdusyntax}[emph={[1]type}]
    type = (*<doctor|master|(bachelor)>*)
  \end{fdusyntax}

  选择论文类型。三种选项分别代表博士学位论文、硕士学位论文和本科
  毕业论文。
\end{function}

\begin{function}{oneside,twoside}
  指明论文的单双面模式，默认为 \opt{twoside}。该选项会影响每章
  的开始位置，还会影响页眉样式。
\end{function}

在双面模式（\opt{twoside}）下，按照通常的排版惯例，每章应只从
奇数页（在右）开始；而在单页模式（\opt{oneside}）下，则可以从
任意页面开始。本模板中，目录、摘要、符号表等均视作章，也按相同
方式排版。

双面模式下，正文部分偶数页（在左）的左页眉显示章标题，奇数页
（在右）的右页眉显示节标题；前置部分的页眉按同样格式显示，但文字
均为对应标题（如“目录”、“摘要”等）。
而在单面模式下，正文部分则页面不分奇偶，均同时显示左、右页眉，
文字分别为章标题和节标题；前置部分只有中间页眉，显示对应标题。

\begin{function}{draft}
  \begin{fdusyntax}[emph={[1]draft}]
    draft = (*<\TFF>*)
  \end{fdusyntax}
  选择是否开启草稿模式，默认关闭。
\end{function}

草稿模式为全局选项，会影响到很多宏包的工作方式。
开启之后，主要的变化有：
\begin{itemize}
  \item 把行溢出的盒子显示为黑色方块；
  \item 不实际插入图片，只输出一个占位方框；
  \item 关闭超链接渲染，也不再生成 PDF 书签；
  \item 显示页面边框。
\end{itemize}

\begin{function}[added=2018-01-31]{config}
  \begin{fdusyntax}[emph={[1]config}]
    config = (*\marg{文件}*)
  \end{fdusyntax}
  用户配置文件的文件名。默认为空，即不载入用户配置文件。
\end{function}

\subsection{参数设置}

\begin{function}{\fdusetup}
  \begin{fdusyntax}[morekeywords={\fdusetup}]
    \fdusetup(*\marg{键值列表}*)
  \end{fdusyntax}
  本模板提供了一系列选项，可由您自行配置。载入文档类之后，以下
  所有选项均可通过统一的命令 \cs{fdusetup} 来设置。
\end{function}

\cs{fdusetup} 的参数是一组由（英文）逗号隔开的选项列表，列表中的
选项通常是 \kvopt{\meta{key}}{\meta{value}} 的形式。部分选项的
\meta{value} 可以省略。对于同一项，后面的设置将会覆盖前面的设置。
在下文的说明中，将用\textbf{粗体}表示默认值。

\cs{fdusetup} 采用 \LaTeX3 风格的键值设置，支持不同类型以及多种
层次的选项设定。键值列表中，“|=|”左右的空格不影响设置；但需注意，
参数列表中不可以出现空行。

与模板选项相同，布尔型的参数可以省略 \kvopt{\meta{选项}}{true}
中的“\kvopt{}{true}”。

另有一些选项包含子选项，如 \opt{style} 和 \opt{info} 等。它们可以
按如下两种等价方式来设定：

\begin{latexexample}[morekeywords={\fdusetup},
    emph={[1]style,cjk-font,font-size,info,title,title*,author,author*,department}]
  \fdusetup{
    style = {cjk-font = adobe, font-size = -4},
    info  = {
      title      = {论动体的电动力学},
      title*     = {On the Electrodynamics of Moving Bodies},
      author     = {阿尔伯特·爱因斯坦},
      author*    = {Albert Einstein},
      department = {物理学系}
    }
  }
\end{latexexample}

或者

\begin{latexexample}[morekeywords={\fdusetup},
    emph={[1]style,cjk-font,font-size,info,title,title*,author,author*,department}]
  \fdusetup{
    style/cjk-font  = adobe,
    style/font-size = -4,
    info/title      = {论动体的电动力学},
    info/title*     = {On the Electrodynamics of Moving Bodies},
    info/author     = {阿尔伯特·爱因斯坦},
    info/author*    = {Albert Einstein},
    info/department = {物理学系}
  }
\end{latexexample}


注意 “|/|” 的前后均不可以出现空白字符。

\subsubsection{论文格式} \label{subsubsec:论文格式}

\begin{function}{style}
  \begin{fdusyntax}[emph={[1]style}]
    style = (*\marg{键值列表}*)
    style/(*\meta{key}*) = (*\meta{value}*)
  \end{fdusyntax}
  该选项包含许多子项目，用于设置论文格式。具体内容见下。
\end{function}


\begin{function}[updated=2019-03-05]{style/font}
  \begin{fdusyntax}[emph={[1]font}]
    font = (*<garamond|libertinus|lm|palatino|(times)|times*|none>*)
  \end{fdusyntax}

  设置西文字体（包括数学字体）。具体配置见表~\ref{tab:font}。
\end{function}

\begin{table}[ht]
\centering
\begin{threeparttable}
  \caption{西文字体配置}
  \label{tab:font}
  \small
  \begin{tabular}{ccccc}
    \toprule
      & \textbf{正文字体} & \textbf{无衬线字体} & \textbf{等宽字体} & \textbf{数学字体} \\
    \midrule
      |garamond|        & EB Garamond         & Libertinus Sans & LM Mono\tnote{a} & Garamond Math   \\
      |libertinus|      & Libertinus Serif    & Libertinus Sans & LM Mono          & Libertinus Math \\
      |lm|              & LM Roman            & LM Sans         & LM Mono          & LM Math         \\
      |palatino|        & TG Pagella\tnote{b} & Libertinus Sans & LM Mono          & TG Pagella Math \\
      |times|           & XITS                & TG Heros        & TG Cursor        & XITS Math       \\
      |times*|\tnote{c} & Times New Roman     & Arial           & Courier New      & XITS Math       \\
    \bottomrule
  \end{tabular}
  \begin{tablenotes}
    \item[a] “LM”是 Latin Modern 的缩写。
    \item[b] “TG”是 TeX Gyre 的缩写。
    \item[c] 本行中，Times New Roman、Arial 和 Courier New 是商业字体，
      不包含在 \TeXLive{} 发行版中，但在 Windows 和 macOS 系统上均默认安装。
  \end{tablenotes}
\end{threeparttable}
\end{table}


\begin{function}[rEXP,updated=2019-03-05]{style/cjk-font}
  \begin{fdusyntax}[emph={[1]cjk-font}]
    cjk-font = (*<adobe|(fandol)|founder|mac|sinotype|sourcehan|windows|none>*)
  \end{fdusyntax}

  设置中文字体。具体配置见表~\ref{tab:cjk-font}。
\end{function}

\begin{table}[ht]
  \caption{中文字体配置}
  \label{tab:cjk-font}
  \centering
  \small
  \begin{tabular}{ccccc}
    \toprule
      & \textbf{正文字体（宋体）} & \textbf{无衬线字体（黑体）} & \textbf{等宽字体（仿宋）} & \textbf{楷体} \\
    \midrule
      |adobe|     & Adobe 宋体      & Adobe  黑体     & Adobe  仿宋  & Adobe 楷体      \\
      |fandol|    & Fandol 宋体     & Fandol 黑体     & Fandol 仿宋  & Fandol 楷体     \\
      |founder|   & 方正书宋        & 方正黑体        & 方正仿宋     & 方正楷体        \\
      |mac|       & （华文）宋体-简 & （华文）黑体-简 & 华文仿宋     & （华文）楷体-简 \\
      |sinotype|  & 华文宋体        & 华文黑体        & 华文仿宋     & 华文楷体        \\
      |sourcehan| & 思源宋体        & 思源黑体        & ---          & ---             \\
      |windows|   & （中易）宋体    & （中易）黑体    & （中易）仿宋 & （中易）楷体    \\
    \bottomrule
  \end{tabular}
\end{table}

启用 \kvopt{font}{none} 或 \kvopt{cjk-font}{none} 之后，模板将关闭
默认西文 / 中文字体设置。此时，您需要自行使用 \cs{setmainfont}、
\cs{setCJKmainfont}、\cs{setmathfont} 等命令来配置字体。


\begin{function}{style/font-size}
  \begin{fdusyntax}[emph={[1]font-size}]
    font-size = (*<(-4)|5>*)
  \end{fdusyntax}

  设置论文的基础字号。
\end{function}


\begin{function}[rEXP,updated=2017-10-14]{style/fullwidth-stop}
  \begin{fdusyntax}[emph={[1]fullwidth-stop}]
    fullwidth-stop = (*<catcode|mapping|(false)>*)
  \end{fdusyntax}

  选择是否把全角实心句点\FSFW 作为默认的句号形状。
  这种句号一般用于科技类文章，以避免与下标“$_o$”或“$_0$”混淆。
\end{function}

选择 \kvopt{fullwidth-stop}{catcode} 或 \opt{mapping} 后，都会实现上述效果。
有所不同的是，在选择 \opt{catcode} 后，只有\emph{显式的}\FSID 会被替换
为\FSFW；但在选择 \opt{mapping} 后，\emph{所有的}\FSID 都会被替换。例如，
如果您用宏保存了一些含有\FSID 的文字，那么在选择 \opt{catcode} 时，其中
的\FSID 不会将被替换为\FSFW。

选项 \kvopt{fullwidth-stop}{mapping} 只在 \XeTeX{} 下有效。使用
\LuaTeX{} 编译时，该选项相当于 \kvopt{fullwidth-stop}{catcode}。

如果您在选择 \kvopt{fullwidth-stop}{mapping} 后仍需要临时显示\FSID，
可以按如下方法操作：
\begin{latexexample}[moretexcs={\CJKfontspec},emph={[1]Mapping}]
  % 请使用 XeTeX 编译
  % 外侧的花括号表示分组
  这是一个句号{\CJKfontspec{(*\meta{字体名}*)}[Mapping=full-stop]。}
\end{latexexample}

\begin{function}{style/footnote-style}
% 这里奇怪的东西是用来控制对齐的。fdusyntax 会吃掉开头的几个
% 空格，因此这里用 X 来占位。
  \begin{fdusyntax}[emph={[1]footnote-style}]
    footnote-style = (*<plain|\\
      XXXX\mbox{}~~~~~~~~~~~~~~~~~libertinus|libertinus*|libertinus-sans|\\
      XXXX\mbox{}~~~~~~~~~~~~~~~~~pifont|pifont*|pifont-sans|pifont-sans*|\\
      XXXX\mbox{}~~~~~~~~~~~~~~~~~xits|xits-sans|xits-sans*>*)
  \end{fdusyntax}
  设置脚注编号样式。西文字体设置会影响其默认取值（见
  表~\ref{tab:footnote-font}）。因此，要使得该选项生效，需将其
  放置在 \opt{font} 选项之后。带有 |sans| 的为相应的无衬线字体
  版本；带有 |*| 的为阴文样式（即黑底白字）。
\end{function}

\begin{table}[ht]
  \caption{西文字体与脚注编号样式默认值的对应关系}
  \label{tab:footnote-font}
  \centering
  \small
  \begin{tabular}{ccccc}
    \toprule
      \textbf{西文字体设置} &
        |libertinus| & |lm|     & |palatino| & |times| \\
    \midrule
      \textbf{脚注编号样式默认值} &
        |libertinus| & |pifont| & |pifont|   & |xits|  \\
    \bottomrule
  \end{tabular}
\end{table}


\begin{function}[added=2017-08-13]{style/hyperlink}
  \begin{fdusyntax}[emph={[1]hyperlink}]
    hyperlink = (*<border|(color)|none>*)
  \end{fdusyntax}

  设置超链接样式。\opt{border} 表示在超链接四周绘制方框；
  \opt{color} 表示用彩色显示超链接；\opt{none} 表示没有特殊装饰，
  可用于生成最终的打印版文稿。
\end{function}

\begin{function}[added=2017-08-13,updated=2021-12-27]{style/hyperlink-color}
  \begin{fdusyntax}[emph={[1]hyperlink-color}]
    hyperlink-color = (*<(default)|classic|material|graylevel|prl>*)
  \end{fdusyntax}
  设置超链接颜色。该选项在 \kvopt{hyperlink}{none} 时无效。
  各选项所代表的颜色见表~\ref{tab:hyperlink-color}。
\end{function}


\begin{table}[ht]
\centering
\small

\newcommand\linkcolorexam[3]{%
  {\small 图~\textcolor[HTML]{#1}{1-2}，
    (\textcolor[HTML]{#1}{3.4})~式} &
  {\small \textcolor[HTML]{#2}{\texttt{https://g.cn}}} &
  {\small 文献~[\textcolor[HTML]{#3}{1}]，
    (\textcolor[HTML]{#3}{Knuth~1986})}}
\begin{threeparttable}
\caption{预定义的超链接颜色方案}
\label{tab:hyperlink-color}
\begin{tabular}{c*{3}{>{\hspace{0.2cm}}c<{\hspace{0.2cm}}}}
  \toprule
    \textbf{选项} & \textbf{链接} & \textbf{URL} & \textbf{引用} \\

  \midrule
    \opt{default}            & \linkcolorexam{990000}{0000B2}{007F00} \\
    \opt{classic}            & \linkcolorexam{FF0000}{0000FF}{00FF00} \\
    \opt{material}\tnote{a}  & \linkcolorexam{E91E63}{009688}{4CAF50} \\
    \opt{graylevel}\tnote{a} & \linkcolorexam{616161}{616161}{616161} \\
    \opt{prl}\tnote{b}       & \linkcolorexam{2D3092}{2D3092}{2D3092} \\
  \bottomrule
\end{tabular}
\begin{tablenotes}

  \item[a] 取自 Material 色彩方案
    （见 \url{https://material.io/guidelines/style/color.html}）。
  \item[b] \textit{Physical Review Letter} 杂志配色。

\end{tablenotes}
\end{threeparttable}
\end{table}



\begin{function}[added=2018-01-25]{style/bib-backend}
  \begin{fdusyntax}[emph={[1]bib-backend}]
    bib-backend = (*<bibtex|biblatex>*)
  \end{fdusyntax}

  选择参考文献的支持方式。选择 \opt{bibtex} 后，将使用 \BibTeX{}
  处理文献，样式由 \pkg{natbib} 宏包负责；选择 \opt{biblatex} 后，
  将使用 \biber{} 处理文献，样式则由 \pkg{biblatex} 宏包负责。
\end{function}

\begin{function}[added=2017-10-28,updated=2018-01-25]{style/bib-style}
  \begin{fdusyntax}[emph={[1]bib-style}]
    bib-style = (*<author-year|(numerical)|\meta{其他样式}>*)
  \end{fdusyntax}
  设置参考文献样式。\opt{author-year} 和 \opt{numerical} 分别对应
  国家标准 GB/T 7714--2015 \cite{gb-t-7714-2015} 中的著者—出版年制和
  顺序编码制。选择 \meta{其他样式} 时，如果 \kvopt{bib-backend}{bibtex}，
  需保证相应的 \file{.bst} 格式文件能被调用；而如果
  \kvopt{bib-backend}{biblatex}，则需保证相应的 \file{.bbx} 格式文件
  能被调用。
\end{function}

\begin{function}[added=2018-01-25]{style/cite-style}
  \begin{fdusyntax}[emph={[1]cite-style}]
    cite-style = (*\marg{引用样式}*)
  \end{fdusyntax}
  选择引用格式。默认为空，即与参考文献样式（著者—出版年制或顺序
  编码制）保持一致。如果手动填写，需保证相应的 \file{.cbx} 格式文件
  能被调用。该选项在 \kvopt{bib-backend}{bibtex} 时无效。
\end{function}

\begin{function}[added=2018-01-25]{style/bib-resource}
  \begin{fdusyntax}[emph={[1]bib-resource}]
    bib-resource = (*\marg{文件}*)
  \end{fdusyntax}
  参考文献数据源。可以是单个文件，也可以是用英文逗号隔开的一组文件。
  如果 \kvopt{bib-backend}{biblatex}，则必须明确给出 \file{.bib}
  后缀名。
\end{function}

\begin{function}[added=2017-08-10]{style/logo}
  \begin{fdusyntax}[emph={[1]logo}]
    logo = (*\marg{文件}*)
  \end{fdusyntax}
  封面中校名图片的文件名。默认值为 \file{fudan-name.pdf}。
\end{function}

\begin{function}[added=2017-08-10]{style/logo-size}
  \begin{fdusyntax}[emph={[1]logo-size}]
    logo-size = (*\marg{宽度}*)
    logo-size = {(*\meta{宽度}*), (*\meta{高度}*)}
  \end{fdusyntax}
  校名图片的大小。默认仅指定了宽度，为 |0.5\textwidth|\/。
  如果仅需指定高度，可在 \meta{宽度} 处填入一个空的分组 |{}|。
\end{function}


\begin{function}[added=2017-07-06]{style/auto-make-cover}
  \begin{fdusyntax}[emph={[1]auto-make-cover}]
    auto-make-cover = (*<\TTF>*)
  \end{fdusyntax}

  是否自动生成论文封面（封一）、指导小组成员名单（封二）和
  声明页（封三）。封面中的各项信息，可通过 \cs{fdusetup} 录入，
  具体请参阅 \ref{subsubsec:信息录入}~节。
\end{function}


\begin{function}[added=2021-09-21]{style/declaration-page}
  \begin{fdusyntax}[emph={[1]declaration-page}]
    declaration-page = (*\marg{文件}*)
  \end{fdusyntax}

  插入扫描版的声明页 PDF 文档。如果为空（默认），则使用模板预定义的声明页。
\end{function}

\begin{function}{\makecoveri,\makecoverii,\makecoveriii}
  用于手动生成论文封面、指导小组成员名单和声明页。这几个命令不能
  确保页码的正确编排，因此除非必要，您应当始终使用自动生成的封面。
\end{function}

\subsubsection{信息录入} \label{subsubsec:信息录入}

\begin{function}{info}
  \begin{fdusyntax}[emph={[1]info}]
    info = (*\marg{键值列表}*)
    info/(*\meta{key}*) = (*\meta{value}*)
  \end{fdusyntax}
  该选项包含许多子项目，用于录入论文信息。具体内容见下。以下带“|*|”
  的项目表示对应的英文字段。
\end{function}


\begin{function}[added=2018-02-01,updated=2019-03-12]{info/degree}
  \begin{fdusyntax}[emph={[1]degree}]
    degree = (*<(academic)|professional>*)
  \end{fdusyntax}

  学位类型，仅适用于博士和硕士学位论文。\opt{academic} 和 \opt{professional}
  分别表示学术学位和专业学位。
\end{function}

\begin{function}{info/title,info/title*}
  \begin{fdusyntax}[emph={[1]title,title*}]
    title  = (*\marg{中文标题}*)
    title* = (*\marg{英文标题}*)
  \end{fdusyntax}
  论文标题。默认会在约 20 个汉字字宽处强制断行，但为了语义的
  连贯以及排版的美观，如果您的标题长于一行，建议使用“|\\|”
  手动断行。
\end{function}

\begin{function}{info/author,info/author*}
  \begin{fdusyntax}[emph={[1]author,author*}]
    author  = (*\marg{姓名}*)
    author* = (*\marg{英文姓名（或拼音）}*)
  \end{fdusyntax}
  作者姓名。
\end{function}

\begin{function}{info/supervisor}
  \begin{fdusyntax}[emph={[1]supervisor}]
    supervisor = (*\marg{姓名}*)
  \end{fdusyntax}
  导师姓名。
\end{function}

\begin{function}{info/department}
  \begin{fdusyntax}[emph={[1]department}]
    department = (*\marg{名称}*)
  \end{fdusyntax}
  院系名称。
\end{function}

\begin{function}{info/major}
  \begin{fdusyntax}[emph={[1]major}]
    major = (*\marg{名称}*)
  \end{fdusyntax}
  专业名称。
\end{function}

\begin{function}{info/student-id}
  \begin{fdusyntax}[emph={[1]student-id}]
    student-id = (*\marg{数字}*)
  \end{fdusyntax}
  作者学号。
\end{function}

复旦大学学号共 11 位，前两位为入学年份，之后一位为学生类型
代码（博士生为 1，硕士生为 2，本科生为 3），接下来的五位为
专业代码，最后三位为顺序号。

\begin{function}{info/school-id}
  \begin{fdusyntax}[emph={[1]school-id}]
    school-id = (*\marg{数字}*)
  \end{fdusyntax}
  学校代码。默认值为 10246（这是复旦大学的学校代码）。
\end{function}

\begin{function}{info/date}
  \begin{fdusyntax}[emph={[1]date}]
    date = (*\marg{日期}*)
  \end{fdusyntax}
  论文完成日期。默认值为文档编译日期（\tn{today}）。
\end{function}


\begin{function}[added=2017-07-04]{info/secret-level}
  \begin{fdusyntax}[emph={[1]secret-level}]
    secret-level = (*<(none)|i|ii|iii>*)
  \end{fdusyntax}

  密级。\opt{i}、\opt{ii}、\opt{iii} 分别表示秘密、机密、绝密；
  \opt{none} 表示论文不涉密，即不显示密级与保密年限。
\end{function}

\begin{function}[added=2017-07-04]{info/secret-year}
  \begin{fdusyntax}[emph={[1]secret-year}]
    secret-year = (*\marg{年限}*)
  \end{fdusyntax}
  保密年限。建议您使用中文，如“五年”。该选项在设置
  \kvopt{secret-level}{none} 时无效。
\end{function}

\begin{function}{info/instructors}
  \begin{fdusyntax}[emph={[1]instructors}]
    instructors = (*\marg{成员 1, 成员 2, ...}*)
  \end{fdusyntax}
  指导小组成员。各成员之间需使用英文逗号隔开。为防止歧义，
  可以用分组括号“|{...}|”把各成员字段括起来。
\end{function}

\begin{function}{info/keywords,info/keywords*}
  \begin{fdusyntax}[emph={[1]keywords,keywords*}]
    keywords  = (*\marg{中文关键字}*)
    keywords* = (*\marg{英文关键字}*)
  \end{fdusyntax}
  关键字列表。各关键字之间需使用英文逗号隔开。为防止歧义，
  可以用分组括号“|{...}|”把各字段括起来。
\end{function}

\begin{function}{info/clc}
  \begin{fdusyntax}[emph={[1]clc}]
    clc = (*\marg{分类号}*)
  \end{fdusyntax}
  中图分类号（CLC）。
\end{function}

\begin{function}[added=2021-09-16]{info/jel}
  \begin{fdusyntax}[emph={[1]jel}]
    jel = (*\marg{分类号}*)
  \end{fdusyntax}
  JEL 分类号，仅适用于部分院系。当被指定时，英文摘要中将仅显示 JEL
  分类号，而不再显示中图分类号。
\end{function}

\subsection{正文编写}

\begin{quotation}
  喬孟符（吉）博學多能，以樂府稱。嘗云：「作樂府亦有法，
  曰\CJKunderdot{鳳頭、豬肚、豹尾}六字是也。」大概起要美麗，中要浩蕩，
  結要響亮。尤貴在首尾貫穿，意思清新。苟能若是，斯可以言樂府矣。
\end{quotation}
\hfill ——陶宗儀《南村輟耕錄·作今樂府法》

\subsubsection{凤头}

\begin{function}{\frontmatter}
  声明前置部分开始。
\end{function}

在本模板中，前置部分包含目录、中英文摘要以及符号表等。
前置部分的页码采用小写罗马字母，并且与正文分开计数。

\begin{function}{\tableofcontents,\listoffigures,\listoftables}
  生成目录。为了生成完整、正确的目录，您至少需要编译\emph{两次}。对于图表
  较多的论文，也可以使用 \cs{listoffigures} 和 \cs{listoftables} 生成单独的
  插图、表格目录。
\end{function}

% TODO: \DescribeEnv{abstract}
% TODO: \DescribeEnv{abstract*}
\begin{function}{abstract}
  \begin{fdusyntax}[emph={[2]abstract}]
    % 中文论文模板 (fduthesis)      % 英文论文模板 (fduthesis-en)
    \begin{abstract}                \begin{abstract}
      (*\meta{中文摘要} \hspace{3.52cm} \meta{英文摘要}*)
    \end{abstract}                  \end{abstract}
  \end{fdusyntax}
\end{function}
\begin{function}[rEXP]{abstract*}
  \begin{fdusyntax}[emph={[2]abstract*}]
    % 中文论文模板 (fduthesis)
    \begin{abstract*}
      (*\meta{英文摘要}*)
    \end{abstract*}
  \end{fdusyntax}
  摘要。中文模板中，不带星号和带星号的版本分别用来输入中文摘要
  和英文摘要；英文模板中没有带星号的版本，您只需输入英文摘要。
\end{function}

摘要的最后，会显示关键字列表以及中图分类号（CLC）或 JEL 分类号。这两项
可通过 \cs{fdusetup} 录入，具体请参阅 \ref{subsubsec:信息录入}~节。

% TODO: \DescribeEnv{notation}
\begin{function}{notation}
  \begin{fdusyntax}[emph={[2]notation}]
    \begin{notation}(*\oarg{列格式说明}*)
      (*\meta{符号 1}*)  &  (*\meta{说明}*)  \\
      (*\meta{符号 2}*)  &  (*\meta{说明}*)  \\
      (*\phantom{\meta{符号 $n$}}*)  (*$\vdots$*)
      (*\meta{符号\ \kern-0.1em$n$}*)  &  (*\meta{说明}*)
    \end{notation}
  \end{fdusyntax}
  符号表。可选参数 \meta{列格式说明}与 \LaTeX{} 中标准表格的列格
  式说明语法一致，默认值为“|lp{7.5cm}|”，即第一列宽度自动调整，
  第二列限宽 \qty{7.5}{cm}，两列均为左对齐。
\end{function}

\subsubsection{猪肚}

\begin{function}{\mainmatter}
  声明主体部分开始。
\end{function}

主体部分是论文的核心，您可以分章节撰写。如有需求，也可以采用
多文件编译的方式。主体部分的页码采用阿拉伯数字。

\begin{function}[updated=2018-01-15]{\footnote}
  \begin{fdusyntax}[deletetexcs={\footnote},morekeywords={\footnote}]
    \footnote(*\marg{脚注文字}*)
  \end{fdusyntax}
  插入脚注。脚注编号样式可利用 \opt{style/footnote-style} 选项控制，
  具体见 \ref{subsubsec:论文格式}~小节。
\end{function}

\begin{function}{\caption}
  \begin{fdusyntax}[deletetexcs={\caption},morekeywords={\caption}]
    \caption(*\marg{图表标题}*)
    \caption(*\oarg{短标题}\marg{长标题}*)
  \end{fdusyntax}
  插入图表标题。可选参数 \meta{短标题} 用于图表目录。在
  \meta{长标题} 中，您可以进行长达多段的叙述；但 \meta{短标题}
  和单独的 \meta{图表标题} 中则不允许分段。
  \cite{刘海洋2013latex入门}
\end{function}

按照排版惯例，建议您将表格的标题放置在绘制表格的命令之前，
而将图片的标题放置在绘图或插图的命令之后。另需注意，
\tn{caption} 命令必须放置在浮动体环境（如 \env{table} 和
\env{figure}）中。

\paragraph{参考文献引用}

\begin{function}{\cite}
  \begin{fdusyntax}[deletetexcs={\cite},morekeywords={\cite}]
    \cite(*\marg{文献标签}*)
    \cite(*\oarg{页码}\marg{文献标签}*)
  \end{fdusyntax}
  插入所引用的文献。可选参数 \meta{页码} 可用来标注引文的页码。在不同的
  参考文献样式中，引用的样式也不尽相同。根据需要，模板还提供了更多的命令
  用来标记引用。顺序编码制和著者—出版年制下的各种引用方式见
  表~\ref{tab:citation-numerical} 和表~\ref{tab:citation-author-year}。
\end{function}

\NewDocumentCommand\verbcite{O{cite}om}{%
  \IfNoValueTF{#2}{%
    \texttt{\textbackslash#1\{#3\}}%
  }{%
    \texttt{\textbackslash#1[#2]\{#3\}}%
  }}

\begin{table}[ht]
  \caption{顺序编码制下的引用样式} \label{tab:citation-numerical}
  \centering
  \small
  \begin{tabularx}{\textwidth}{cCll}
    \toprule
      \textbf{引用方式} &
      \textbf{排版效果} &
      \textbf{\kvopt{bib-backend}{bibtex}} &
      \textbf{\kvopt{bib-backend}{biblatex}} \\
    \midrule
    单篇文献 &
      引文\cite{texbook} &
      \verbcite{texbook} &
      同右 \\
    多篇文献 &
      引文\cite{texbook,companion} &
      \verbcite{texbook,companion} &
      同右 \\
    标注页码 &
      引文\cite[126--137]{texbook} &
      \verbcite[cite][126--137]{texbook} &
      同右 \\
    标注作者 &
      \citet{texbook}指出 &
      \verbcite[citet]{texbook} &
      \verbcite[authornumcite]{texbook} \\
    标注作者、页码 &
      \citet[42]{texbook}指出 &
      \verbcite[citet][42]{texbook} &
      \verbcite[authornumcite][42]{texbook} \\
    非上标 &
      引文{\citestyle{numbers}\cite{texbook}} &
      \verbcite[parencite]{texbook} &
      同右 \\
    \bottomrule
  \end{tabularx}
\end{table}

\begin{table}[ht]
  \caption{著者—出版年制下的引用样式} \label{tab:citation-author-year}
  \centering
  \small
  \def\!{\kern-1.5pt}
  \begin{tabularx}{\textwidth}{cCll}
    \toprule
      \textbf{引用方式} &
      \textbf{排版效果} &
      \textbf{\kvopt{bib-backend}{bibtex}} &
      \textbf{\kvopt{bib-backend}{biblatex}} \\
    \midrule
    单篇文献 &
      (\citeauthor{texbook}\!, \citeyear{texbook}\!) &
      \verbcite[citep]{texbook} &
      \verbcite{texbook} \\
    多篇文献 &
      (\citeauthor{texbook}\!, \citeyear{texbook}\!;
        \citeauthor{companion}\!, \citeyear{companion}\!) &
      \verbcite[citep]{texbook,companion} &
      \verbcite{texbook,companion} \\
    标注页码 &
      (\citeauthor{texbook}\!, \citeyear{texbook}\!)%
        \textsuperscript{126--137} &
      \verbcite[citep][126--137]{texbook} &
      \verbcite[cite][126--137]{texbook} \\
    标注作者 &
      \citeauthor{texbook}(\citeyear{texbook}\!) &
      \verbcite[citet]{texbook} &
      同右 \\
    标注作者、页码 &
      \citeauthor{texbook}(\citeyear{texbook}\!)\textsuperscript{42} &
      \verbcite[citet][42]{texbook} &
      同右 \\
    \bottomrule
  \end{tabularx}
\end{table}

\paragraph{定理类环境}

% TODO: \DescribeEnv{proof}
\begin{function}{axiom,corollary,definition,example,lemma,
  proof,theorem}
  \begin{fdusyntax}[emph={[2]proof}]
    \begin{proof}(*\oarg{小标题}*)
      (*\meta{证明过程}*)
    \end{proof}
  \end{fdusyntax}
  一系列预定义的数学环境。具体含义见表~\ref{tab:theorem}。
\end{function}

\begin{table}[ht]
  \caption{预定义的数学环境} \label{tab:theorem}
  \centering
  \small
  \begin{tabular}{*{8}{c}}
    \toprule
      \textbf{名称} &
        \env{axiom} & \env{corollary} & \env{definition} & \env{example} &
        \env{lemma} & \env{proof}     & \env{theorem} \\
    \midrule
      \textbf{含义} &
        公理 & 推论 & 定义 & 例 &
        引理 & 证明 & 定理 \\
    \bottomrule
  \end{tabular}
\end{table}

证明环境（\env{proof}）的最后会添加证毕符号“$\QED$”。要确保
该符号在正确的位置显示，您需要按照 \ref{subsec:编译方式}~节
中的有关说明编译\emph{两次}。

\begin{function}[updated=2017-12-12]{\newtheorem}
  \begin{fdusyntax}[deletetexcs={\newtheorem},
      morekeywords={\newtheorem,\newtheorem*}]
    \newtheorem(*\oarg{选项}\marg{环境名}\marg{标题}*)
    \newtheorem*(*\oarg{选项}\marg{环境名}\marg{标题}*)
    \begin(*\marg{环境名}\oarg{小标题}*)
      (*\meta{内容}*)
    \end(*\marg{环境名}*)
  \end{fdusyntax}
  声明新的定理类环境（数学环境）。带星号的版本表示不进行编号，
  并且会默认添加证毕符号“$\QED$”。声明后，即可同预定义的数学环境
  一样使用。
\end{function}

事实上，表~\ref{tab:theorem} 中预定义的环境正是通过以下方式定义的：
\begin{latexexample}[deletetexcs={\newtheorem},
    morekeywords={\newtheorem,\newtheorem*}]
  \newtheorem*{proof}{证明}
  \newtheorem{axiom}{公理}
  \newtheorem{corollary}{定理}
  ...
\end{latexexample}

与 \cs{fdusetup} 相同，\cs{newtheorem} 的可选参数 \meta{选项}
也为一组键值列表。可用的选项见下。注意您无需输入“|theorem/|”。

\begin{function}{theorem/style}
  \begin{fdusyntax}[emph={[1]style}]
    style = (*<(plain)|margin|change|\\
      XXXXXX\mbox{}~~~~~~~~break|marginbreak|changebreak>*)
  \end{fdusyntax}
  定理类环境的总体样式。
\end{function}

\begin{function}{theorem/header-font}
  \begin{fdusyntax}[emph={[1]header-font}]
    header-font = (*\marg{字体}*)
  \end{fdusyntax}
  定理头（即标题）的字体。中文模板默认为 \tn{sffamily}，即无衬线体
  （黑体）；英文模板默认为 |\bfseries\upshape|，即加粗直立体。
\end{function}

\begin{function}{theorem/body-font}
  \begin{fdusyntax}[emph={[1]body-font}]
    body-font = (*\marg{字体}*)
  \end{fdusyntax}
  定理内容的字体。中文模板默认为 \tn{fdu@kai}，即楷体；英文模板
  默认为 \tn{itshape}，即斜体。
\end{function}

\begin{function}{theorem/qed}
  \begin{fdusyntax}[emph={[1]qed}]
    qed = (*\marg{符号}*)
  \end{fdusyntax}
  定理结束标记（即证毕符号）。如果用 \cs{newtheorem} 声明定理，则默认为空；
  用 \cs{newtheorem*} 声明，则默认为 |\ensuremath{\QED}|，即“$\QED$”。
\end{function}

\begin{function}{theorem/counter}
  \begin{fdusyntax}[emph={[1]counter}]
    counter = (*\marg{计数器}*)
  \end{fdusyntax}
  定理计数器，表示定理编号在 \meta{计数器} 的下一级，并会随 \meta{计数器}
  的变化而清零。\cite{刘海洋2013latex入门}默认为 |chapter|，表示按章编号。
  使用 \cs{newtheorem*} 时，该选项无效。
\end{function}

\subsubsection{豹尾}

\begin{function}{\backmatter}
  声明后置部分开始。
\end{function}

后置部分包含参考文献、声明页等。

\begin{function}[updated=2018-01-25]{\printbibliography}
  \begin{fdusyntax}[morekeywords={\printbibliography}]
    \printbibliography(*\oarg{选项}*)
  \end{fdusyntax}
  打印参考文献列表。如果 \kvopt{bib-backend}{bibtex}，则 \meta{选项}
  无效，相当于 \tn{bibliography} \texttt{\marg{文献数据库}}，其中的
  \meta{文献数据库} 可利用 \opt{style/bib-resource} 选项指定，具体见
  \ref{subsubsec:论文格式}~小节；而如果 \kvopt{bib-backend}{biblatex}，
  则该命令由 \pkg{biblatex} 宏包直接提供，可用选项请参阅其文档
  \cite{biblatex}。
\end{function}

\section{宏包依赖情况}

使用不同编译方式、指定不同选项，会导致宏包依赖情况有所不同。
具体如下：
\begin{itemize}
  \item 在任何情况下，本模板都会\emph{显式}调用以下宏包
    （或文档类）：
    \begin{itemize}
      \item \pkg{xtemplate} 和 \pkg{l3keys2e}，用于扩展 \LaTeX3 编程环境。
        它们属于 \pkg{l3packages} 宏集。
      \item \cls{ctexbook}，提供中文排版的通用框架。属于 \CTeX{}
        宏集 \cite{CTeX}。
      \item \pkg{amsmath}，对 \LaTeX{} 的数学排版功能进行了
        全面扩展。属于 \AmSLaTeX{} 套件。
      \item \pkg{unicode-math}，负责处理 Unicode 编码的
        OpenType 数学字体。
      \item \pkg{geometry}，用于调整页面尺寸。
      \item \pkg{fancyhdr}，处理页眉页脚。
      \item \pkg{footmisc}，处理脚注。
      \item \pkg{ntheorem}，提供增强版的定理类环境。
      \item \pkg{graphicx}，提供图形插入的接口。
      \item \pkg{longtable}，长表格（允许跨页）支持。
      \item \pkg{caption}，用于设置题注。
      \item \pkg{xcolor}，提供彩色支持。
      \item \pkg{hyperref}，提供交叉引用、超链接、电子书签等功能。
    \end{itemize}
  \item 开启 \kvopt{style/footnote-style}{pifont} 后，会调用
    \pkg{pifont} 宏包。它属于 \pkg{psnfss} 套件。
  \item 开启 \kvopt{style/bib-backend}{bibtex} 后，会调用
    \pkg{natbib} 宏包，并依赖 \BibTeX{} 程序。参考文献样式由
    \pkg{gbt7714} 宏包提供 \cite{natbib,gbt7714}。
  \item 开启 \kvopt{style/bib-backend}{biblatex} 后，会调用
    \pkg{biblatex} 宏包，并依赖 \biber{} 程序。参考文献样式由
    \pkg{biblatex-gb7714-2015} 宏包提供
    \cite{biblatex,biblatex-gb7714-2015}。
\end{itemize}

这里只列出了本模板直接调用的宏包。这些宏包自身的调用情况，
此处不再具体展开。如有需要，请参阅相关文档。

\begin{thebibliography}{99}

\newcommand\urlprefix{\newline\hspace*{\fill}}
\let\OldUrl=\url
\renewcommand\url[2][]{{\small\textit{#1}~\OldUrl{#2}}}
\newcommand\CTANurl[2][]{{%
  \small\textit{#1}~\href{http://mirror.ctan.org/#2}{\texttt{CTAN://#2}}}}

\subsection{图书}

\bibitem[Knuth(1986)]{texbook}
\textsc{Knuth D E}.
\newblock \textit{The \TeX book: Computers \& Typesetting, volume A} [M].
\newblock Boston: Addison--Wesley Publishing Company, 1986
\urlprefix \CTANurl[源代码%
  \footnote{此代码只可作为学习之用。未经 Knuth 本人同意，您不应当编译此文档。}：]%
  {systems/knuth/dist/tex/texbook.tex}

\bibitem[Mittelbach et~al.(2004)]{companion}
\textsc{Mittelbach F} and \textsc{Goossens M}.
\newblock \textit{The \LaTeX{} Companion} [M].
\newblock 2nd ed.
\newblock Boston: Addison--Wesley Publishing Company, 2004

\bibitem[()]{胡伟2017latex2e}
胡伟.
\newblock \textit{\LaTeXe{} 文类和宏包学习手册} [M].
\newblock 北京: 清华大学出版社, 2017

\bibitem[()]{刘海洋2013latex入门}
刘海洋.
\newblock \textit{\LaTeX{} 入门} [M].
\newblock 北京: 电子工业出版社, 2013

\subsection{标准、规范}

\bibitem[()]{gb-t-7713.1-2006}
国务院学位委员会办公室, 全国信息与文献标准化技术委员会.
\newblock \textit{学位论文编写规则: GB/T 7713.1--2006} [S].
\newblock 北京: 中国标准出版社, 2007

\bibitem[()]{gb-t-7714-2015}
全国信息与文献标准化技术委员会.
\newblock \textit{信息与文献\quad 参考文献著录规则: GB/T 7714--2015} [S].
\newblock 北京: 中国标准出版社, 2015

\bibitem[()]{gb-t-15834-2011}
教育部语言文字信息管理司.
\newblock \textit{标点符号用法: GB/T 15834--2011} [S/OL].
\newblock 北京: 中国标准出版社, 2012
\urlprefix\url{http://www.moe.gov.cn/ewebeditor/uploadfile/2015/01/13/20150113091548267.pdf}

\bibitem[()]{clreq}
W3C.
\newblock \textit{中文排版需求（Requirements for Chinese Text Layout）} [EB/OL].
\newblock (2020-08-17)
\urlprefix\url{https://w3c.github.io/clreq/}

\bibitem[()]{复旦大学论文规范}
复旦大学图书馆, 复旦大学研究生院.
\newblock \textit{复旦大学博士、硕士学位论文规范} [EB/OL].
\newblock 2017 年 3 月修订版.
\newblock (2017-03-27)
\urlprefix\url{http://www.gs.fudan.edu.cn/_upload/article/4c/a8/a82545ef443b9c057c14ba13782c/c883c6f3-6d7f-410c-8f30-d8bde6fcb990.doc}

\subsection{宏包、模版}

\bibitem[()]{source2e}
\textsc{Braams J}, \textsc{Carlisle D}, \textsc{Jeffrey A}, et al.
\newblock \textit{The \LaTeXe{} Sources} [CP/OL].
\newblock (2020-02-02)
\urlprefix\url{https://ctan.org/pkg/latex}
\urlprefix\CTANurl[源代码：]{macros/latex/base/source2e.pdf}

\bibitem[()]{CTeX}
\textsc{CTEX.ORG}.
\newblock \textit{\CTeX{} 宏集手册} [EB/OL].
\newblock version 2.5.4,
\newblock (2020-08-23)
\urlprefix\url{https://ctan.org/pkg/ctex}
\urlprefix\CTANurl[文档及源代码：]{language/chinese/ctex/ctex.pdf}

\bibitem[()]{xeCJK}
\textsc{CTEX.ORG}.
\newblock \textit{\pkg{xeCJK} 宏包} [EB/OL].
\newblock version 3.8.5,
\newblock (2020-08-23)
\urlprefix\url{https://ctan.org/pkg/xecjk}
\urlprefix\CTANurl[文档及源代码：]{macros/xetex/latex/xecjk/xeCJK.pdf}

\bibitem[()]{natbib}
\textsc{Daly P W}.
\newblock \textit{Natural Sciences Citations and References} [EB/OL].
\newblock version 8.31b,
\newblock (2010-09-13)
\urlprefix\url{https://ctan.org/pkg/natbib}
\urlprefix\CTANurl[文档及源代码：]{macros/latex/contrib/natbib/natbib.pdf}

\bibitem[()]{source3}
\textsc{The \LaTeX3 Project}.
\newblock \textit{The \LaTeX3 Sources} [CP/OL].
\newblock (2020-08-07)
\urlprefix\url{https://ctan.org/pkg/l3kernel}
\urlprefix\CTANurl[源代码：]{macros/latex/contrib/l3kernel/source3.pdf}

\bibitem[()]{biblatex}
\textsc{Lehman P}, \textsc{Kime P}, \textsc{Boruvka A}, et al.
\newblock \textit{The \pkg{biblatex} Package} [EB/OL].
\newblock version 3.15a,
\newblock (2020-08-23)
\urlprefix\url{https://ctan.org/pkg/biblatex}
\urlprefix\CTANurl[文档：]{macros/latex/contrib/biblatex/doc/biblatex.pdf}

\bibitem[()]{lshort}
\textsc{Oetiker T}, \textsc{Partl H}, \textsc{Hyna I}, et al.
\newblock \textit{The Not So Short Introduction to \LaTeXe{}: Or \LaTeXe{} in 139 minutes} [EB/OL].
\newblock version 6.2,
\newblock (2018-02-28)
\urlprefix\url{https://ctan.org/pkg/lshort-english}
\urlprefix\CTANurl[文档：]{info/lshort/english/lshort.pdf}

\bibitem[()]{lshort-zh-cn}
\textsc{Oetiker T}, \textsc{Partl H}, \textsc{Hyna I}, et al.
\newblock \textit{一份（不太）简短的 \LaTeXe{} 介绍: 或 112 分钟了解 \LaTeXe{}} [EB/OL].
\newblock \CTeX{} 开发小组, 译.
\newblock 原版版本 version 6.2, 中文版本 version 6.02,
\newblock (2020-08-03)
\urlprefix\url{https://ctan.org/pkg/lshort-zh-cn}
\urlprefix\CTANurl[文档：]{info/lshort/chinese/lshort-zh-cn.pdf}

\bibitem[()]{biblatex-gb7714-2015}
胡振震.
\newblock \textit{符合 GB/T 7714-2015 标准的 biblatex 参考文献样式} [EB/OL].
\newblock version 1.0v,
\newblock (2020-07-21)
\urlprefix\url{https://ctan.org/pkg/biblatex-gb7714-2015}
\urlprefix\CTANurl[文档：]{biblatex-contrib/biblatex-gb7714-2015/biblatex-gb7714-2015.pdf}

\bibitem[()]{gbt7714}
李泽平（\textsc{Zeping L}）.
\newblock \textit{GB/T 7714-2015 \BibTeX{} Style} [EB/OL].
\newblock version 2.0.2,
\newblock (2020-06-08)
\urlprefix\url{https://ctan.org/pkg/gbt7714}
\urlprefix\CTANurl[文档：]{biblio/bibtex/contrib/gbt7714/gbt7714.pdf}

\bibitem[()]{cquthesis}
李振楠.
\newblock \textit{\textsc{CquThesis}：重庆大学毕业论文 \LaTeX{} 模板} [EB/OL].
\newblock version 1.50,
\newblock (2020-06-23)
\urlprefix\url{https://ctan.org/pkg/cquthesis}
\urlprefix\CTANurl[文档及源代码：]{macros/latex/contrib/cquthesis/cquthesis.pdf}

\bibitem[()]{thuthesis}
清华大学 TUNA 协会.
\newblock \textit{\textsc{ThuThesis}：清华大学学位论文模板} [EB/OL].
\newblock version 6.1.3,
\newblock (2020-07-09)
\urlprefix\url{https://ctan.org/pkg/thuthesis}
\urlprefix\CTANurl[文档及源代码：]{macros/latex/contrib/thuthesis/thuthesis.pdf}

\emph{以下模版未收录至 CTAN，但仍然保持活跃更新。}

\bibitem[()]{sjtuthesis}
\textsc{SJTUG}.
\newblock \textit{上海交通大学学位论文模板} [EB/OL].
\newblock version 1.0.0rc7,
\newblock (2020-07-31)
\urlprefix\url{https://github.com/sjtug/SJTUThesis}
\urlprefix\url[源代码仓库：]{https://github.com/sjtug/SJTUTeX}

\bibitem[()]{ustcthesis}
\textsc{USTC \TeX{} User Group}.
\newblock \textit{中国科学技术大学学位论文 \LaTeX{} 模板} [EB/OL].
\newblock version 3.2.1,
\newblock (2020-04-03)
\urlprefix\url{https://github.com/ustctug/ustcthesis}

\bibitem[()]{ucasthesis}
\textsc{mohuangrui}.
\newblock \textit{\pkg{ucasthesis} 国科大学位论文 \LaTeX{} 模板} [EB/OL].
\newblock (2020-07-28)
\urlprefix\url{https://github.com/mohuangrui/ucasthesis}

\emph{以下模版现已停止更新。}

\bibitem[()]{pandoxie2014fduthesislatex}
\textsc{Pandoxie}.
\newblock \textit{Fudan University-Latex Template} [EB/OL].
\newblock (2014-06-07)
\urlprefix\url{https://github.com/Pandoxie/FDU-Thesis-Latex}

\bibitem[()]{richard2016fudanthesis}
\textsc{richard}.
\newblock \textit{复旦大学硕士学位论文模板} [EB/OL].
\newblock (2016-01-31)
\urlprefix\url{https://github.com/richarddzh/fudan-thesis}

\bibitem[()]{数院毕业论文格式}
复旦大学数学科学学院.
\newblock \textit{毕业论文格式 tex 版和 word 版} [EB/OL].
\urlprefix\url{http://math.fudan.edu.cn/show.aspx?info_lb=664&flag=101&info_id=1816}

\bibitem[()]{数院毕业论文格式更新}
复旦大学数学科学学院.
\newblock \textit{毕业论文格式: Word、\TeX{} 模板更新} [EB/OL].
\urlprefix\url{http://math.fudan.edu.cn/Show.aspx?info_lb=664&info_id=1855&flag=101}

\subsection{其他}

\bibitem[()]{wright2009dtxfile}
\textsc{Wright J}.
\newblock \textit{A model dtx file} [EB/OL].
\newblock (2009-10-06)
\urlprefix\url{https://www.texdev.net/2009/10/06/a-model-dtx-file/}

\bibitem[()]{孔雀计划}
刘庆（\textsc{Eric Q L}）.
\newblock \textit{孔雀计划：中文字体排印的思路} [EB/OL].
\urlprefix\url{https://thetype.com/kongque/}

\end{thebibliography}

\clearpage

\end{document}